\section{Kết luận}
\label{sec:conclusion}

Trong bài viết này, lần đầu tiên chúng tôi thiết lập rõ ràng mối liên hệ giữa các phương pháp phát hiện bất thường truyền thống dựa trên sự phụ thuộc và các phương pháp phát hiện bất thường theo ngữ cảnh. Trên cơ sở này, chúng tôi đề xuất một cách tiếp cận mới để phát hiện và giải thích sự bất thường theo ngữ cảnh. Cụ thể, chúng tôi sử dụng Rừng hồi quy phân vị để phát triển phương pháp phát hiện bất thường chính xác và dễ hiểu, QCAD, nhằm khám phá sự phụ thuộc giữa các tính năng. QCAD có thể xử lý các tập dữ liệu dạng bảng với các đặc điểm ngữ cảnh hỗn hợp và các đặc trưng hành vi số. Kết quả thử nghiệm mở rộng trên nhiều bộ dữ liệu tổng hợp và thực tế khác nhau chứng minh rằng QCAD vượt trội hơn các phương pháp phát hiện bất thường tiên tiến nhất trong việc xác định các bất thường theo ngữ cảnh về độ chính xác và khả năng giải thích.

Từ nghiên cứu điển hình về dữ liệu cầu thủ bóng đá, chúng tôi kết luận rằng QCAD có thể phát hiện những bất thường theo ngữ cảnh có ý nghĩa và hữu ích mà chuyên gia miền có thể giải thích trực tiếp. Các hình ảnh trực quan hóa dựa trên Beanplot giúp giải thích lý do tại sao một đối tượng nhất định (không) được coi là bất thường trong bối cảnh của nó. Điều này rất quan trọng vì việc phát hiện sự bất thường là một vấn đề không được giám sát và những lời giải thích có thể giúp các nhà phân tích và chuyên gia miền xác minh kết quả. Nó cũng mở ra cơ hội phát hiện sự bất thường do con người hướng dẫn, trong đó phản hồi từ nhà phân tích có thể được sử dụng để hướng dẫn phân tích.

Trong tương lai, do QCAD chỉ có thể xử lý các tính năng tĩnh nên chúng tôi dự định mở rộng QCAD sang cài đặt phát trực tuyến.

\section*{Tuyên bố và Tuyên bố}



\bmhead{Phê duyệt đạo đức}Dữ liệu về con người và động vật liên quan đến nghiên cứu này được cung cấp công khai và do đó không cần phải phê duyệt.

\bmhead{Tài trợ}
Công việc này được hỗ trợ bởi Dự án 4 của chương trình nghiên cứu Digital Twin, một chương trình Phối cảnh TTW với số dự án P18-03 được tài trợ chủ yếu bởi Hội đồng Nghiên cứu Hà Lan (NWO). Tất cả các ý kiến, phát hiện, kết luận và khuyến nghị trong bài viết này là của các tác giả và không nhất thiết phản ánh quan điểm của các cơ quan tài trợ.

\bmhead{Xung đột lợi ích} (Các) tác giả tuyên bố không có xung đột lợi ích tiềm ẩn với
tôn trọng việc nghiên cứu, quyền tác giả và/hoặc việc xuất bản cuốn sách này
bài viết.

\bmhead{Có sẵn dữ liệu và tài liệu}
Để có thể tái tạo, toàn bộ mã và bộ dữ liệu được cung cấp tại \url{https://github.com/ZhongLIFR/QCAD}.

\bmhead{Đóng góp về quyền tác giả}

\textit{Zhong Li}: Khái niệm hóa, Phương pháp luận, Xác nhận, Điều tra, Viết, Trực quan hóa, Quản lý Dự án. \textit{Matthijs van Leeuwen}: Phương pháp luận, Xác thực, Viết, Nhận tài trợ.

\begin{thebibliography}{66}
\providecommand{\natexlab}[1]{#1}
\providecommand{\url}[1]{{#1}}
\providecommand{\urlprefix}{URL }
\providecommand{\doi}[1]{\url{https://doi.org/#1}}
\providecommand{\eprint}[2][]{\url{#2}}
 \bibcommenthead

\bibitem[{Aggarwal and Sathe(2017)}]{aggarwal2017outlier}
Aggarwal CC, Sathe S (2017) Outlier ensembles: An introduction. Springer

\bibitem[{Ahmad et~al(2017)Ahmad, Munir, Bhatti, Aftab, and
  Raza}]{ahmad2017survival}
Ahmad T, Munir A, Bhatti SH, et~al (2017) Survival analysis of heart failure
  patients: A case study. PloS one 12(7):e0181,001

\bibitem[{Ahmed et~al(2016{\natexlab{a}})Ahmed, Mahmood, and
  Hu}]{ahmed2016survey1}
Ahmed M, Mahmood AN, Hu J (2016{\natexlab{a}}) A survey of network anomaly
  detection techniques. Journal of Network and Computer Applications 60:19--31

\bibitem[{Ahmed et~al(2016{\natexlab{b}})Ahmed, Mahmood, and
  Islam}]{ahmed2016survey2}
Ahmed M, Mahmood AN, Islam MR (2016{\natexlab{b}}) A survey of anomaly
  detection techniques in financial domain. Future Generation Computer Systems
  55:278--288

\bibitem[{Angiulli and Pizzuti(2002)}]{angiulli2002fast}
Angiulli F, Pizzuti C (2002) Fast outlier detection in high dimensional spaces.
  In: European conference on principles of data mining and knowledge discovery,
  Springer, pp 15--27

\bibitem[{Babbar and Chawla(2012)}]{babbar2012mining}
Babbar S, Chawla S (2012) Mining causal outliers using gaussian bayesian
  networks. In: 2012 IEEE 24th International Conference on Tools with
  Artificial Intelligence, IEEE, pp 97--104

\bibitem[{Breiman(2001)}]{breiman2001random}
Breiman L (2001) Random forests. Machine learning 45:5--32

\bibitem[{Breiman et~al(2017)Breiman, Friedman, Olshen, and
  Stone}]{breiman2017classification}
Breiman L, Friedman JH, Olshen RA, et~al (2017) Classification and regression
  trees. Routledge

\bibitem[{Breunig et~al(2000)Breunig, Kriegel, Ng, and Sander}]{breunig2000lof}
Breunig MM, Kriegel HP, Ng RT, et~al (2000) Lof: identifying density-based
  local outliers. In: Proceedings of the 2000 ACM SIGMOD international
  conference on Management of data, pp 93--104

\bibitem[{Buczak and Guven(2015)}]{buczak2015survey}
Buczak AL, Guven E (2015) A survey of data mining and machine learning methods
  for cyber security intrusion detection. IEEE Communications surveys \&
  tutorials 18(2):1153--1176

\bibitem[{Cabero et~al(2021)Cabero, Epifanio, Pi{\'e}rola, and
  Ballester}]{cabero2021archetype}
Cabero I, Epifanio I, Pi{\'e}rola A, et~al (2021) Archetype analysis: A new
  subspace outlier detection approach. Knowledge-Based Systems 217:106,830

\bibitem[{Cai et~al(2013)Cai, He, and Man}]{cai2013spatial}
Cai Q, He H, Man H (2013) Spatial outlier detection based on iterative
  self-organizing learning model. Neurocomputing 117:161--172

\bibitem[{Calikus et~al(2021)Calikus, Nowaczyk, Bouguelia, and
  Dikmen}]{calikus2021wisdom}
Calikus E, Nowaczyk S, Bouguelia MR, et~al (2021) Wisdom of the contexts:
  Active ensemble learning for contextual anomaly detection. arXiv preprint
  arXiv:210111560

\bibitem[{Campos et~al(2016)Campos, Zimek, Sander, Campello, Micenkov{\'a},
  Schubert, Assent, and Houle}]{campos2016evaluation}
Campos GO, Zimek A, Sander J, et~al (2016) On the evaluation of unsupervised
  outlier detection: measures, datasets, and an empirical study. Data mining
  and knowledge discovery 30(4):891--927

\bibitem[{Chandola et~al(2009)Chandola, Banerjee, and
  Kumar}]{chandola2009anomaly}
Chandola V, Banerjee A, Kumar V (2009) Anomaly detection: A survey. ACM
  computing surveys (CSUR) 41(3):1--58

\bibitem[{F{\"a}rber et~al(2010)F{\"a}rber, G{\"u}nnemann, Kriegel, Kr{\"o}ger,
  M{\"u}ller, Schubert, Seidl, and Zimek}]{farber2010using}
F{\"a}rber I, G{\"u}nnemann S, Kriegel HP, et~al (2010) On using class-labels
  in evaluation of clusterings. In: MultiClust: 1st international workshop on
  discovering, summarizing and using multiple clusterings held in conjunction
  with KDD, p~1

\bibitem[{Fokkema et~al(2022)Fokkema, de~Heide, and van
  Erven}]{fokkema2022attribution}
Fokkema H, de~Heide R, van Erven T (2022) Attribution-based explanations that
  provide recourse cannot be robust. arXiv preprint arXiv:220515834

\bibitem[{Friedman(1937)}]{friedman1937use}
Friedman M (1937) The use of ranks to avoid the assumption of normality
  implicit in the analysis of variance. Journal of the american statistical
  association 32(200):675--701

\bibitem[{Goldstein and Dengel(2012)}]{goldstein2012histogram}
Goldstein M, Dengel A (2012) Histogram-based outlier score (hbos): A fast
  unsupervised anomaly detection algorithm. KI-2012: Poster and Demo Track pp
  59--63

\bibitem[{Gower(1971)}]{gower1971general}
Gower JC (1971) A general coefficient of similarity and some of its properties.
  Biometrics pp 857--871

\bibitem[{Harrison~Jr and Rubinfeld(1978)}]{harrison1978hedonic}
Harrison~Jr D, Rubinfeld DL (1978) Hedonic housing prices and the demand for
  clean air. Journal of environmental economics and management 5(1):81--102

\bibitem[{Hawkins(1980)}]{hawkins1980identification}
Hawkins DM (1980) Identification of outliers, vol~11. Springer

\bibitem[{Hayes and Capretz(2014)}]{hayes2014contextual}
Hayes MA, Capretz MA (2014) Contextual anomaly detection in big sensor data.
  In: 2014 IEEE International Congress on Big Data, IEEE, pp 64--71

\bibitem[{Hong and Hauskrecht(2015)}]{hong2015multivariate}
Hong C, Hauskrecht M (2015) Multivariate conditional anomaly detection and its
  clinical application. In: Proceedings of the AAAI Conference on Artificial
  Intelligence

\bibitem[{Huang et~al(2003)Huang, Fan, Lee, and Yu}]{huang2003cross}
Huang Ya, Fan W, Lee W, et~al (2003) Cross-feature analysis for detecting
  ad-hoc routing anomalies. In: 23rd International Conference on Distributed
  Computing Systems, 2003. Proceedings., IEEE, pp 478--487

\bibitem[{Hwang et~al(2009)Hwang, Kim, Kim, and Seah}]{hwang2009survey}
Hwang I, Kim S, Kim Y, et~al (2009) A survey of fault detection, isolation, and
  reconfiguration methods. IEEE transactions on control systems technology
  18(3):636--653

\bibitem[{Kampstra(2008)}]{kampstra2008beanplot}
Kampstra P (2008) Beanplot: A boxplot alternative for visual comparison of
  distributions. Journal of statistical software 28(1):1--9

\bibitem[{Kandanaarachchi et~al(2020)Kandanaarachchi, Mu{\~n}oz, Hyndman, and
  Smith-Miles}]{kandanaarachchi2020normalization}
Kandanaarachchi S, Mu{\~n}oz MA, Hyndman RJ, et~al (2020) On normalization and
  algorithm selection for unsupervised outlier detection. Data Mining and
  Knowledge Discovery 34(2):309--354

\bibitem[{Koenker and Hallock(2001)}]{koenker2001quantile}
Koenker R, Hallock KF (2001) Quantile regression. Journal of economic
  perspectives 15(4):143--156

\bibitem[{Kriegel et~al(2009)Kriegel, Kr{\"o}ger, Schubert, and
  Zimek}]{kriegel2009outlier}
Kriegel HP, Kr{\"o}ger P, Schubert E, et~al (2009) Outlier detection in
  axis-parallel subspaces of high dimensional data. In: Pacific-asia conference
  on knowledge discovery and data mining, Springer, pp 831--838

\bibitem[{Kriegel et~al(2012)Kriegel, Kr{\"o}ger, Schubert, and
  Zimek}]{kriegel2012outlier}
Kriegel HP, Kr{\"o}ger P, Schubert E, et~al (2012) Outlier detection in
  arbitrarily oriented subspaces. In: 2012 IEEE 12th international conference
  on data mining, IEEE, pp 379--388

\bibitem[{Kuo et~al(2018)Kuo, Li, and Kifer}]{kuo2018detecting}
Kuo YH, Li Z, Kifer D (2018) Detecting outliers in data with correlated
  measures. In: Proceedings of the 27th ACM International Conference on
  Information and Knowledge Management, pp 287--296

\bibitem[{Lei et~al(2018)Lei, G’Sell, Rinaldo, Tibshirani, and
  Wasserman}]{lei2018distribution}
Lei J, G’Sell M, Rinaldo A, et~al (2018) Distribution-free predictive
  inference for regression. Journal of the American Statistical Association
  113(523):1094--1111

\bibitem[{Li et~al(2022)Li, Zhu, and van Leeuwen}]{li2022survey}
Li Z, Zhu Y, van Leeuwen M (2022) A survey on explainable anomaly detection.
  arXiv preprint arXiv:221006959

\bibitem[{Liang and Parthasarathy(2016)}]{liang2016robust}
Liang J, Parthasarathy S (2016) Robust contextual outlier detection: Where
  context meets sparsity. In: Proceedings of the 25th ACM International on
  Conference on Information and Knowledge Management, pp 2167--2172

\bibitem[{Liu et~al(2008)Liu, Ting, and Zhou}]{liu2008isolation}
Liu FT, Ting KM, Zhou ZH (2008) Isolation forest. In: 2008 eighth ieee
  international conference on data mining, IEEE, pp 413--422

\bibitem[{Liu et~al(2018)Liu, Shin, and Hu}]{liu2018contextual}
Liu N, Shin D, Hu X (2018) Contextual outlier interpretation. In: Proceedings
  of the 27th International Joint Conference on Artificial Intelligence, pp
  2461--2467

\bibitem[{Lu et~al(2020{\natexlab{a}})Lu, Liu, Li, Le, and
  Liu}]{lu2020dependency}
Lu S, Liu L, Li J, et~al (2020{\natexlab{a}}) Dependency-based anomaly
  detection: Framework, methods and benchmark. arXiv preprint arXiv:201106716

\bibitem[{Lu et~al(2020{\natexlab{b}})Lu, Liu, Li, Le, and Liu}]{lu2020lopad}
Lu S, Liu L, Li J, et~al (2020{\natexlab{b}}) Lopad: A local prediction
  approach to anomaly detection. Advances in Knowledge Discovery and Data
  Mining 12085:660

\bibitem[{Lundberg and Lee(2017)}]{lundberg2017unified}
Lundberg SM, Lee SI (2017) A unified approach to interpreting model
  predictions. Advances in neural information processing systems 30

\bibitem[{Meghanath et~al(2018)Meghanath, Pai, and
  Akoglu}]{meghanath2018conout}
Meghanath M, Pai D, Akoglu L (2018) Conout: Con textual outlier detection with
  multiple contexts: Application to ad fraud. In: Joint European Conference on
  Machine Learning and Knowledge Discovery in Databases, Springer, pp 139--156

\bibitem[{Meinshausen(2006)}]{meinshausen2006quantile}
Meinshausen N (2006) Quantile regression forests. Journal of Machine Learning
  Research 7(6):983--999

\bibitem[{Micenkov{\'a} et~al(2014)Micenkov{\'a}, McWilliams, and
  Assent}]{micenkova2014learning}
Micenkov{\'a} B, McWilliams B, Assent I (2014) Learning outlier ensembles: The
  best of both worlds--supervised and unsupervised. In: Proceedings of the ACM
  SIGKDD 2014 Workshop on Outlier Detection and Description under Data
  Diversity (ODD2). New York, NY, USA, Citeseer, pp 51--54

\bibitem[{Micenkov{\'a} et~al(2015)Micenkov{\'a}, McWilliams, and
  Assent}]{micenkova2015learning}
Micenkov{\'a} B, McWilliams B, Assent I (2015) Learning representations for
  outlier detection on a budget. arXiv preprint arXiv:150708104

\bibitem[{Nemenyi(1963)}]{nemenyi1963distribution}
Nemenyi PB (1963) Distribution-free multiple comparisons. Princeton University

\bibitem[{Nguyen et~al(2013)Nguyen, M{\"u}ller, Vreeken, Keller, and
  B{\"o}hm}]{nguyen2013cmi}
Nguyen HV, M{\"u}ller E, Vreeken J, et~al (2013) Cmi: An information-theoretic
  contrast measure for enhancing subspace cluster and outlier detection. In:
  Proceedings of the 2013 SIAM International Conference on Data Mining, SIAM,
  pp 198--206

\bibitem[{Noto et~al(2010)Noto, Brodley, and Slonim}]{noto2010anomaly}
Noto K, Brodley C, Slonim D (2010) Anomaly detection using an ensemble of
  feature models. In: 2010 ieee international conference on data mining, IEEE,
  pp 953--958

\bibitem[{Pang et~al(2021)Pang, Shen, Cao, and Hengel}]{pang2021deep}
Pang G, Shen C, Cao L, et~al (2021) Deep learning for anomaly detection: A
  review. ACM Computing Surveys (CSUR) 54(2):1--38

\bibitem[{Panjei et~al(2022)Panjei, Gruenwald, Leal, Nguyen, and
  Silvia}]{panjei2022survey}
Panjei E, Gruenwald L, Leal E, et~al (2022) A survey on outlier explanations.
  The VLDB Journal 31(5):977--1008

\bibitem[{Pasillas-D{\'\i}az and Ratt{\'e}(2016)}]{pasillas2016unsupervised}
Pasillas-D{\'\i}az JR, Ratt{\'e} S (2016) An unsupervised approach for
  combining scores of outlier detection techniques, based on similarity
  measures. Electronic notes in theoretical computer science 329:61--77

\bibitem[{Salvador et~al(2004)Salvador, Chan, and
  Brodie}]{salvador2004learning}
Salvador S, Chan P, Brodie J (2004) Learning states and rules for time series
  anomaly detection. In: FLAIRS conference, pp 306--311

\bibitem[{Scutari et~al(2019)Scutari, Scutari, and MMPC}]{scutari2019package}
Scutari M, Scutari MM, MMPC HP (2019) Package ‘bnlearn’. Bayesian network
  structure learning, parameter learning and inference, R package version 4(1)

\bibitem[{Segal and Xiao(2011)}]{segal2011multivariate}
Segal M, Xiao Y (2011) Multivariate random forests. Wiley interdisciplinary
  reviews: Data mining and knowledge discovery 1(1):80--87

\bibitem[{Seger(2018)}]{seger2018investigation}
Seger C (2018) An investigation of categorical variable encoding techniques in
  machine learning: binary versus one-hot and feature hashing

\bibitem[{Smets et~al(2009)Smets, Verdonk, and Jordaan}]{smets2009discovering}
Smets K, Verdonk B, Jordaan EM (2009) Discovering novelty in spatio/temporal
  data using one-class support vector machines. In: 2009 International Joint
  Conference on Neural Networks, IEEE, pp 2956--2963

\bibitem[{Song et~al(2007)Song, Wu, Jermaine, and Ranka}]{song2007conditional}
Song X, Wu M, Jermaine C, et~al (2007) Conditional anomaly detection. IEEE
  Transactions on knowledge and Data Engineering 19(5):631--645

\bibitem[{Spinosa and Carvalho(2005)}]{spinosa2005support}
Spinosa EJ, Carvalho A (2005) Support vector machines for novel class detection
  in bioinformatics. Genet Mol Res 4(3):608--15

\bibitem[{Tang et~al(2015)Tang, Pei, Bailey, and Dong}]{tang2015mining}
Tang G, Pei J, Bailey J, et~al (2015) Mining multidimensional contextual
  outliers from categorical relational data. Intelligent Data Analysis
  19(5):1171--1192

\bibitem[{Teng(1999)}]{teng1999correcting}
Teng CM (1999) Correcting noisy data. In: ICML, Citeseer, pp 239--248

\bibitem[{Valko et~al(2011)Valko, Kveton, Valizadegan, Cooper, and
  Hauskrecht}]{valko2011conditional}
Valko M, Kveton B, Valizadegan H, et~al (2011) Conditional anomaly detection
  with soft harmonic functions. In: 2011 IEEE 11th International Conference on
  Data Mining, IEEE, pp 735--743

\bibitem[{Wang et~al(2019)Wang, Bah, and Hammad}]{wang2019progress}
Wang H, Bah MJ, Hammad M (2019) Progress in outlier detection techniques: A
  survey. Ieee Access 7:107,964--108,000

\bibitem[{Wong et~al(2003)Wong, Moore, Cooper, and Wagner}]{wong2003bayesian}
Wong WK, Moore AW, Cooper GF, et~al (2003) Bayesian network anomaly pattern
  detection for disease outbreaks. In: Proceedings of the 20th International
  Conference on Machine Learning (ICML-03), pp 808--815

\bibitem[{Xu et~al(2021)Xu, Wang, Jian, Huang, Wang, Liu, and
  Li}]{xu2021beyond}
Xu H, Wang Y, Jian S, et~al (2021) Beyond outlier detection: Outlier
  interpretation by attention-guided triplet deviation network. In: Proceedings
  of the Web Conference 2021, pp 1328--1339

\bibitem[{Yaramakala and Margaritis(2005)}]{yaramakala2005speculative}
Yaramakala S, Margaritis D (2005) Speculative markov blanket discovery for
  optimal feature selection. In: Fifth IEEE International Conference on Data
  Mining (ICDM'05), IEEE, pp 4--pp

\bibitem[{Zhao et~al(2019)Zhao, Nasrullah, and Li}]{zhao2019pyod}
Zhao Y, Nasrullah Z, Li Z (2019) Pyod: A python toolbox for scalable outlier
  detection. Journal of Machine Learning Research 20:1--7

\bibitem[{Zheng et~al(2017)Zheng, Brantley, Lauvaux, and
  Li}]{zheng2017contextual}
Zheng G, Brantley SL, Lauvaux T, et~al (2017) Contextual spatial outlier
  detection with metric learning. In: Proceedings of the 23rd ACM SIGKDD
  International Conference on Knowledge Discovery and Data Mining, pp
  2161--2170

\end{thebibliography}



\appendix
